\section{Configuration System Overview}
The centralized configuration system is a key architectural component of this face recognition project. It addresses the common challenge of managing numerous parameters and settings across different modules without hardcoding values directly in the code. This approach significantly improves maintainability, flexibility, and the ability to experiment with different settings.

\section{Design Principles}
The configuration system was designed with the following principles in mind:
\begin{itemize}
    \item \textbf{Centralization:} All configuration parameters are defined in a single location
    \item \textbf{Separation of Concerns:} Code logic is separated from configuration values
    \item \textbf{Hierarchical Organization:} Settings are logically grouped by component
    \item \textbf{Default Values:} Sensible defaults ensure the system works out of the box
    \item \textbf{Runtime Configuration:} Settings can be overridden via command-line arguments
    \item \textbf{Environment Independence:} Different environments can use different configuration files
\end{itemize}

\section{Implementation Details}
The configuration system consists of two main components:
\begin{enumerate}
    \item A YAML-based configuration file (\texttt{config.yaml})
    \item A settings module (\texttt{src/config/settings.py})
\end{enumerate}

\subsection{Configuration File Structure}
The configuration file uses a hierarchical structure organized by functional components:

\begin{lstlisting}[language=yaml, caption=Configuration File Structure]
# App-level settings
app:
  default_mode: "camera"
  verbose: false
  
# Camera settings
camera:
  default_index: 0
  show_gui: true
  resolution:
    width: 640
    height: 480
  raspi_ip: "192.168.1.17"

# Detection model settings
detection:
  default_model: "mediapipe"
  mediapipe:
    min_detection_confidence: 0.5
  yolov8:
    model_path: "models/yolov8n-face.pt"
    confidence_threshold: 0.5
    iou_threshold: 0.45
  yolov8_onnx:
    model_path: "models/yolov8n_face.onnx"
    confidence_threshold: 0.5
    iou_threshold: 0.45
    input_shape: [640, 640]
    
# Recognition model settings
recognition:
  model_path: "models/facenet.tflite"
  input_shape: [160, 160]
  normalize_embeddings: true
  
# ... additional component settings ...
\end{lstlisting}

\subsection{Settings Module}
The settings module is responsible for loading the configuration file and providing convenient access to settings:

\begin{lstlisting}[language=Python, caption=Settings Module Implementation]
import os
import yaml
from typing import Dict, Any, Union

# Load the configuration file
with open(os.path.join(os.path.dirname(__file__), "config.yaml"), "r") as f:
    settings = yaml.safe_load(f)

# Convert nested dictionaries to dot-accessible objects
class DotDict(dict):
    __getattr__ = dict.get
    __setattr__ = dict.__setitem__
    __delattr__ = dict.__delitem__
    
    def __init__(self, dct):
        for key, value in dct.items():
            if isinstance(value, dict):
                value = DotDict(value)
            self[key] = value

# Convert settings to dot notation for easier access
settings = DotDict(settings)

# Extract component-specific settings for direct import
app_settings = settings.get("app", {})
camera_settings = settings.get("camera", {})
detection_settings = settings.get("detection", {})
recognition_settings = settings.get("recognition", {})
verification_settings = settings.get("verification", {})
alignment_settings = settings.get("alignment", {})
drowsiness_settings = settings.get("drowsiness", {})
database_settings = settings.get("database", {})
api_settings = settings.get("api", {})

# Helper function to get file paths
def get_path(path_key: str) -> str:
    """Get a file path from the configuration.
    
    Args:
        path_key: Dot-separated path key (e.g., "database.user_embeddings_file")
    
    Returns:
        The configured path
    """
    components = path_key.split(".")
    current = settings
    for comp in components:
        current = current.get(comp, {})
    
    if not isinstance(current, str):
        raise ValueError(f"Path {path_key} not found in configuration")
    
    return current
\end{lstlisting}

\section{Usage Patterns}
The configuration system can be used in various ways throughout the application:

\subsection{Direct Component Settings Access}
The most common usage pattern is to import component-specific settings:

\begin{lstlisting}[language=Python]
# Import specific component settings
from src.config import detection_settings

# Use in a class or function
class Detector:
    def __init__(self):
        model_type = detection_settings.get("default_model", "mediapipe")
        # ... initialization code ...
\end{lstlisting}

\subsection{Accessing Nested Settings}
For deeper nested settings, the dot notation can be used:

\begin{lstlisting}[language=Python]
# Import the main settings object
from src.config import settings

# Access nested settings
resolution_width = settings.camera.resolution.width
\end{lstlisting}

\subsection{Getting File Paths}
The \texttt{get\_path} helper function simplifies accessing file paths:

\begin{lstlisting}[language=Python]
from src.config import get_path

# Get a database file path
db_path = get_path("database.user_embeddings_file")
\end{lstlisting}

\section{Overriding Configuration}
The configuration system supports overriding settings via command-line arguments. The main application entry point (\texttt{run.py}) parses command-line arguments and uses them to override default configurations:

\begin{lstlisting}[language=Python]
# Command-line argument parsing in run.py
parser.add_argument("--detector", type=str, 
                    default=detection_settings.get("default_model", "mediapipe"),
                    choices=["mediapipe", "yolov8", "yolov8_onnx"],
                    help="Face detector model to use")
\end{lstlisting}

This allows users to experiment with different settings without modifying the configuration file.

\section{Benefits and Impact}
The implementation of this centralized configuration system has yielded several benefits:

\begin{itemize}
    \item \textbf{Reduced Code Duplication:} Parameters are defined once and reused throughout the application
    \item \textbf{Improved Maintainability:} Changes to settings can be made in one place
    \item \textbf{Enhanced Experimentation:} Different parameter combinations can be tested without code changes
    \item \textbf{Environmental Adaptability:} Different configuration files can be used for development and production
    \item \textbf{Cleaner Code:} Logic is separated from configuration values, improving readability
\end{itemize}

This configuration system has significantly simplified the development, testing, and deployment of the face recognition application, allowing for rapid experimentation with different models and parameters.